% Generated draft: requires mathematical and source review.
\documentclass[11pt]{article}
\usepackage[margin=25mm]{geometry}
\usepackage{fontspec}
\setmainfont{Latin Modern Roman}
\usepackage{amsmath,amssymb,mathtools,unicode-math}
\setmathfont{Latin Modern Math}
\usepackage{enumitem,needspace,xurl,hyperref}
\setlength{\parindent}{0pt}
\setlength{\parskip}{5pt}
\setlength{\emergencystretch}{3em}
\newcommand{\N}{\mathbb{N}}
\newcommand{\Z}{\mathbb{Z}}
\newcommand{\Q}{\mathbb{Q}}
\newcommand{\R}{\mathbb{R}}
\newcommand{\C}{\mathbb{C}}
\newcommand{\F}{\mathbb{F}}
\newcommand{\A}{\mathbb{A}}
\newcommand{\PP}{\mathbb{P}}
\newcommand{\E}{\mathbb{E}}
\newcommand{\eps}{\varepsilon}
\newcommand{\dd}{\,\mathrm{d}}
\begin{document}
\section*{UnsolvedMath ID: 10600056}
\textbf{Source problem code:} catalogue entry
\subsection{Definitions and mathematical statement}
Let $\mathcal G_{\mathrm{rig}}^{\mathrm{or}}$ denote the source's class of rigid graph-links with orientable atoms. Neither the rigid graph-link move system nor the atom construction is defined in the supplied excerpt; these cannot be replaced without justification by ordinary links, arbitrary graphs, or all virtual links.

Problem-list item 56 asks to construct generalizations of the Frobenius-extension framework and of the Rasmussen invariant for $\mathcal G_{\mathrm{rig}}^{\mathrm{or}}$. In a common algebraic convention, a commutative Frobenius $R$-algebra is a finite projective $R$-algebra $A$ equipped with an $R$-linear trace $\epsilon:A\to R$ such that
\[
 A\longrightarrow\operatorname{Hom}_R(A,R),\qquad
 a\longmapsto\bigl(b\longmapsto\epsilon(ab)\bigr)
\]
is an isomorphism. Such algebraic data underlie Frobenius constructions of link homology; this definition is context, not a choice of the particular extension or deformation requested by the source.

The intended research goal is to specify corresponding algebraic/homological data and a Rasmussen-type invariant that are well defined under the rigid graph-link equivalences, retaining the orientable-atom restriction. The word ``Rasmussen'' is read as referring to the invariant, as in the registry summary, but the coefficient field, deformation, filtration, normalization, and expected recovery or functoriality properties are not supplied. In particular, no graph-link analogue of a slice-genus bound is prescribed by this excerpt.
\subsection{Short English statement}
Extend Frobenius-based link homology and a Rasmussen-type invariant to rigid graph-links with orientable atoms. The rigid move system and the exact algebraic and filtered constructions remain to be specified.
\subsection{Sources}
\begin{itemize}
\item Local registry: \url{data/unsolvedmath/problems.json}, permanent UnsolvedMath ID $10600056$, exact source code \texttt{catalogue entry}.
\item \emph{Unsolved Problems in Virtual Knot Theory and Combinatorial Knot Theory}, revised survey, problem-list item 56: \url{https://arxiv.org/abs/1409.2823}. Supplied access date: 2026-07-29; extraction: source-tex; both components of the source item are retained.
\item Supplied related titles: ``Khovanov homology of graph-links,'' \texttt{arXiv:1005.xxxx} (2010), and ``On the virtual Rasmussen invariant,'' \texttt{arXiv:1603.xxxx} (2016). These identifiers are incomplete placeholders, not verified references.
\end{itemize}
\end{document}
